\scanpage{57}
\uline{N.B.} a) Ceci ne caractérise pas les classes algébriques elles-$=$s, mais seulement \unsure{le|les} $\mathbb Q_\ell$-\unsure{espace|espaces} \unsure{engendré|engendrés}.

b) De $=$, on ne peut caractériser ainsi l'image de $\pi$, mais seulement l'algèbre enveloppante\ldots.

Passant à la limite sur les extensions finies de $K$, on trouve
\[
\left\{\begin{aligned}
&H^{2i}(\bar X,\mathbb Q_\ell(i))^{\unsure{\mathfrak g|g}_\ell(\bar K/K)}
\xleftarrow{\sim}\unsure{\mathscr A|A|\mathcal A}^{i}(\bar X)\otimes_{\mathbb Q}\mathbb Q_\ell\\
&\unsure{\mathscr U|\mathscr W|\mathcal U}^i(X)\otimes_{\mathbb Q}\mathbb Q_\ell\quad\text{et l'algèbre enveloppante de}\\
&\hspace{8em}\mathfrak g_\ell(\bar K/K)\ \text{ds }\operatorname{End}_{\mathbb Q_\ell}(H^i(\bar X,\mathbb Q_\ell(0))\\
&\hspace{8em}\text{sont commutants l'un}\\
&\hspace{8em}\text{de l'autre.}
\end{aligned}\right.
\]

\uline{Question} Comment caractériser l'image de $\mathfrak g_\ell(\bar K/K)$ dans $\operatorname{End}_{\mathbb Q_\ell}(H^i(\bar X,\mathbb Q_\ell(0))$ ??

\unsure{Cela|Cel} doit être une sous Algèbre de Lie \uline{réductive}, on pt la considérer comme plongée ds $H^{2n-i}(\bar X)\otimes H^i(\bar X)(n)$.
% archive sheet number: 56

